\section{Competitive Equilibrium}
\label{sec:rewrite-competition}

I now replace the monopolistic developer with competitive AI suppliers.
This comparison separates the incentive to improve AI from the use of
technology already developed. I focus on $\sigma>1$, where both
labor-bottleneck and AI-dominated growth are possible. I first characterize
the competitive equilibrium and then compare it with the monopoly
equilibria in Section~\ref{sec:rewrite-regimes}.

\subsection{Characterization}
\label{subsec:rewrite-competitive-equilibrium}

Consider now the case where multiple AI suppliers have free access to the same efficiency $B$. Entry is free,
and any improvement becomes immediately available to every supplier,
without royalties or other compensation to its developer. There are no
research subsidies or publicly financed research. Suppliers and households
take prices and the common technology as given. Preferences, final-good
production, and the RSI technology remain unchanged.

Competition sets $p_X=1/B$. Since $X=BU$, revenue then exactly covers
inference compute: $p_XX=U$. Final-good firms continue to solve
problem~\eqref{eq:rewrite-final-firm-problem}, with production given by
equations~\eqref{eq:rewrite-output}--\eqref{eq:rewrite-ces}.
Their demand for AI services in equation~\eqref{eq:rewrite-ai-price}
then implies $U=(1-\alpha)s_XY$, since $X=B_0U$. The household continues to choose consumption and saving according to
the Euler equation~\eqref{eq:rewrite-euler} and the transversality
condition~\eqref{eq:rewrite-household-tvc}. Capital accumulation follows
the resource constraint~\eqref{eq:rewrite-resource} with $M=0$.
Competitive suppliers' optimality replaces
the monopolist's pricing, research, and shadow-price conditions.

As labor's income share vanishes, the net return on capital under
marginal-cost pricing approaches
\begin{equation}
 \mathcal R^{\mathrm{mc}}(B)
 =\alpha\left[(1-\alpha)B\,
 \omega_X^{\sigma/(\sigma-1)}\right]^{(1-\alpha)/\alpha}-\delta.
 \label{eq:rewrite-marginal-cost-return}
\end{equation}
Let $\mathcal B^{\mathrm{mc}}(\sigma)$ be the unique efficiency at which
$\mathcal R^{\mathrm{mc}}(\mathcal B^{\mathrm{mc}}(\sigma))=\rho+\gamma$.
Unlike the monopoly economy, the competitive economy cannot improve
from $B_0$ toward $\overline B$. Its long-run regime depends on the
efficiency inherited at the initial date.

\Needspace{24\baselineskip}
\begin{proposition}[Competitive equilibrium without private research]
\label{prop:rewrite-competitive-equilibrium}
Under the free-access arrangement described above, suppose $\sigma>1$.
For every $A_0,N_0,K_0>0$ and $0<B_0<\overline B$, there is a unique competitive equilibrium. At every date,
\begin{equation}
 M=0,\qquad B=B_0,\qquad p_{X}=1/B_0,\qquad \Pi=0.
 \label{eq:rewrite-competitive-no-research}
\end{equation}
Its long-run outcomes are as follows:
\begin{enumerate}[label=(\roman*),leftmargin=*,itemsep=0.5em]
\item If $B_0<\mathcal B^{\mathrm{mc}}(\sigma)$, labor remains a bottleneck:
\[
 g_{Y/N}\to\gamma,\qquad g_w\to\gamma,\qquad
 r\to\rho+\gamma,\qquad \lim_{t\to\infty}wL/Y>0.
\]
\item If $B_0>\mathcal B^{\mathrm{mc}}(\sigma)$, growth is AI-dominated:
\begin{gather}
 g_{Y/N}\to\mathcal R^{\mathrm{mc}}(B_0)-\rho>\gamma,\qquad
 g_w\to\gamma+
 \frac{\mathcal R^{\mathrm{mc}}(B_0)-\rho-\gamma}{\sigma}>\gamma,
 \label{eq:rewrite-competitive-ai-growth}\\
 r\to\mathcal R^{\mathrm{mc}}(B_0),\qquad wL/Y\to0.\nonumber
\end{gather}
\item At $B_0=\mathcal B^{\mathrm{mc}}(\sigma)$, both growth rates
converge to $\gamma$ and $r\to\rho+\gamma$, but $wL/Y\to0$.
\end{enumerate}
\end{proposition}
\begin{proof}
See Appendix~\ref{app:rewrite-proofs},
\hyperref[proof:rewrite-competitive-equilibrium]{Proof of
Proposition~\ref*{prop:rewrite-competitive-equilibrium}}.
\end{proof}

The absence of RSI does not eliminate AI-dominated growth. At fixed
efficiency, $X=B_0U$ can continue expanding through greater compute use.
When $B_0$ exceeds the competitive threshold
$\mathcal B^{\mathrm{mc}}(\sigma)$, capital and AI services
grow together faster than effective labor. This is the capital-accumulation
mechanism identified in Section~\ref{subsec:rewrite-equilibrium-existence},
now sustained without private research. Below the threshold, diminishing
returns bring the interest rate back to $\rho+\gamma$. Since $B=B_0$,
the competitive economy cannot escape this bottleneck through its own
research, even if the upper bound would permit much higher efficiency.
The equality case
is a boundary: capital per unit of effective labor grows without bound,
but its growth rate tends to zero, so the labor share vanishes without a
permanent growth premium.

AI-industry revenue remains distinct from profit. In the competitive
AI-dominated regime, $p_XX/Y\to1-\alpha$, just as under monopoly.
Here, however, all of that revenue pays for inference compute, and net
AI-industry profit is zero. A large AI revenue share therefore need not
represent monopoly rents.

\subsection{Comparison with monopoly}
\label{subsec:rewrite-monopoly-threshold}

I first compare the efficiency required under the two pricing rules,
holding the technology fixed. This isolates the effect of pricing from
the difference in research incentives.

\Needspace{12\baselineskip}
\begin{proposition}[Monopoly pricing and the AI-efficiency threshold]
\label{prop:rewrite-monopoly-threshold}
Suppose AI and labor are substitutes $(\sigma>1)$. Let
$\mathcal B(\sigma)$ and $\mathcal B^{\mathrm{mc}}(\sigma)$ denote the
AI-efficiency thresholds under monopoly pricing and marginal-cost
pricing, respectively.
These thresholds identify the efficiency needed for the return on
capital to reach $\rho+\gamma$ in the limit where labor's income share
vanishes. Monopoly pricing requires a higher AI efficiency:
\begin{equation}
 \mathcal B^{\mathrm{mc}}(\sigma)
 =(1-\alpha)\mathcal B(\sigma)<\mathcal B(\sigma).
 \label{eq:rewrite-marginal-cost-threshold}
\end{equation}
\end{proposition}
\begin{proof}
See Appendix~\ref{app:rewrite-proofs},
\hyperref[proof:rewrite-monopoly-threshold]{Proof of
Proposition~\ref*{prop:rewrite-monopoly-threshold}}.
\end{proof}

To understand why monopoly pricing raises the threshold, I compare the
same AI technology with the same capital and effective labor under the
two pricing rules. At marginal cost, AI services are cheaper, so final-good
firms use more of them and produce more with the same capital. This raises
the output--capital ratio and therefore the return on capital. For the
threshold comparison, I evaluate both returns as labor's income share
tends to zero. In this limit, monopoly pricing requires a higher AI
efficiency for the return on capital to reach the benchmark $\rho+\gamma$.
With $\alpha=0.33$, marginal-cost pricing reaches this benchmark at only
$67.0\%$ of the AI efficiency required under monopoly pricing.

Starting from the same initial state, the two economies need not retain
the same technology. Competitive efficiency remains at $B_0$, whereas
efficiency approaches $\overline B$ along the monopoly equilibria
characterized in Proposition~\ref{prop:rewrite-equilibrium-regimes}.
For example, if
\begin{equation}
 \mathcal B^{\mathrm{mc}}(\sigma)<B_0<\overline B<\mathcal B(\sigma),
 \label{eq:rewrite-competition-escapes-bottleneck}
\end{equation}
the competitive economy has AI-dominated growth, while any monopoly
equilibrium characterized by
Proposition~\ref{prop:rewrite-capped-substitutes-labor-existence}
remains labor-bottlenecked. This comparison applies to a nonempty set of
common initial states: choose
$\overline B\in(\mathcal B^{\mathrm{mc}}(\sigma),\mathcal B(\sigma))$
and $B_0$ sufficiently close to it in the monopoly existence construction.
The competitive equilibrium exists from each of those same states.

The efficiency gain from research must also be considered.
Conditional on both economies converging to AI-dominated
regimes, their limiting returns satisfy
\begin{equation}
 \mathcal R^{\mathrm{mc}}(B_0)\gtrless\mathcal R(\overline B)
 \quad\Longleftrightarrow\quad
 B_0\gtrless(1-\alpha)\overline B.
 \label{eq:rewrite-competitive-monopoly-growth-comparison}
\end{equation}
Equality holds on one side if and only if it holds on the other.
The same ordering applies to long-run output-per-worker and wage growth.
Thus, when both equilibria are AI-dominated, monopoly delivers faster
long-run growth if $B_0<(1-\alpha)\overline B$; competition delivers
faster growth if the inequality is reversed. The research-induced
efficiency gain must therefore be large enough to offset the monopoly
markup. This comparison is conditional on the existence of both
equilibrium trajectories from the same initial state. The competitive
existence result is global, but the monopoly construction in
Proposition~\ref{prop:rewrite-equilibrium-regimes} is local in initial
stocks. Equation~\eqref{eq:rewrite-competitive-monopoly-growth-comparison}
does not extend that existence result to initial states outside those
neighborhoods.

\subsection{Research incentives and policy}
\label{subsec:rewrite-competition-policy}

Holding $K$, $AL$, and $B$ fixed, marginal-cost pricing maximizes
$Y-X/B$ and leaves more resources available for consumption or capital
accumulation. In the free-access economy, however, it also removes the
private return to research. The comparison does not rank welfare:
that requires comparing consumption throughout the transition, including
the resource cost and future benefits of research under monopoly.

This raises a classic time-inconsistency problem \citep{kydlandprescott1977}.
A government may promise monopoly protection to encourage research, but
later wish to remove the markup once the technology has been developed.
Anticipating this policy reversal can weaken the developer's incentive to
invest in research. In the present setting, continued RSI makes the trade-off relevant
even after a valuable AI technology already exists: opening access to
that technology changes the rewards for improving it further.

An announced rule that ends monopoly protection at a specified efficiency
level is different from reneging on a promise. It nevertheless changes
the developer's problem. Without compensation, reaching that level would
remove future monopoly rents, potentially giving the developer an
incentive to delay reaching it. The threshold
$\mathcal B^{\mathrm{mc}}(\sigma)$ identifies the AI efficiency at which the
limiting return on capital under marginal-cost pricing reaches $\rho+\gamma$.
It does not identify the optimal time to end monopoly protection.
Choosing that time requires accounting for the
research response and comparing consumption paths under alternative rules.

\Needspace{3\baselineskip}
Patent buyouts offer a possible way to reward innovation without retaining
the monopoly markup. \citet{kremer1998} studies government purchases of
patents that are generally placed in the public domain. Applied here,
the idea would be to compensate the developer for the rights to an
existing AI technology and allow competing suppliers to use it.
Competitive supply at the same efficiency would require access to the
algorithms and model weights.
The buyout would also need a credible valuation and payment rule, a
source of financing, and arrangements to reward subsequent improvements.
Incentives for further development depend on
whether developers can appropriate its benefits. These issues are
particularly important for RSI, where today's technology is also an
input into tomorrow's research. I leave the design of such a policy,
including its effects on research and welfare, to an extension of the model.
